%=====================================================================
\chapter{Selection and Order Statistics}\label{ch:selection}
%=====================================================================
\locator{3}{Part III --- Sorting, Selection and the Structures Beneath}

\begin{outcomes}
By the end of this chapter you will be able to:
\begin{tight}
  \item \textbf{Define} the $i$-th order statistic and \textbf{distinguish}
        selection from sorting.
  \item \textbf{Derive} quickselect's $\bigTh{n}$ expected and $\bigTh{n^{2}}$
        worst-case running times from the recurrences, and \textbf{explain} why
        recursing on one side changes the master-theorem case.
  \item \textbf{Describe} the median-of-medians algorithm and \textbf{state}
        why it gives a worst-case linear bound.
  \item \textbf{Judge} when the deterministic algorithm is worth its constant.
  \item \textbf{Predict} how selection's cost varies with $k$.
\end{tight}
\end{outcomes}

\begin{prereq}
\chref{quicksort} --- quickselect is quicksort with one recursive call deleted,
and the randomisation argument is the same one. \chref{master} for the
recurrence.
\end{prereq}

\begin{keyterms}
order statistic $\bullet$ selection $\bullet$ median $\bullet$ quickselect
$\bullet$ prune and search $\bullet$ median of medians $\bullet$ introselect
$\bullet$ worst-case linear $\bullet$ tail call
\end{keyterms}

\section{Asking for Less}

\begin{definitionbox}[Order Statistic, Selection]
The \term{$i$-th order statistic} of a set of $n$ elements is the $i$-th
smallest. The 1st is the minimum, the $n$-th the maximum, and the
$\lfloor (n+1)/2 \rfloor$-th the \term{median}.

\smallskip
The \term{selection problem} is: given an array and an index $i$, return the
$i$-th order statistic.
\end{definitionbox}

\begin{conceptbox}[Sorting answers more than was asked]
Selection is obviously solvable in $\bigTh{n\lg n}$: sort, then index. The
question is whether that is wasteful, and it plainly is --- sorting also
determines the relative order of every \emph{other} pair, which nobody
requested.

\smallskip
\chref{lowerbound} proved that \emph{sorting} cannot be done in less than
$\lg(n!)$ comparisons. That theorem says nothing about selection, because
selection has only $n$ possible answers rather than $n!$, so its decision tree
needs only $\lg n$ leaves' worth of height from the counting argument --- a
useless bound, but an honest sign that the problem is easier.

\smallskip
And it is: selection is $\bigTh{n}$, in the worst case, deterministically.
Doing \emph{less} than sorting turns out to cost asymptotically less.
\end{conceptbox}

\section{Quickselect}

\dsfig{\aoaalg{\algname{Quickselect}$(A, \mathit{lo}, \mathit{hi}, k)$}{
\While{\textbf{true}}
  \If{$\mathit{hi} - \mathit{lo} = 1$}
    \State \Return $A[\mathit{lo}]$
  \EndIf
  \State $p \gets$ \algname{Random-Partition}$(A, \mathit{lo}, \mathit{hi})$
  \If{$p = k$}
    \State \Return $A[p]$
      \Comment{the pivot \emph{is} the answer}
  \ElsIf{$k < p$}
    \State $\mathit{hi} \gets p$
      \Comment{discard the right side entirely}
  \Else
    \State $\mathit{lo} \gets p + 1$
      \Comment{discard the left side entirely}
  \EndIf
\EndWhile
}}{\algname{Quickselect}. Identical to \algname{Quicksort} except that one of
the two recursive calls is \emph{deleted} rather than made. Written as a loop
because the remaining call is in tail position --- which also removes
\chref{quicksort}'s stack-overflow failure.}{quickselect}

\begin{theorem}[Quickselect's expected time]\label{thm:qsel}
With a uniformly random pivot, \algname{Quickselect} runs in $\bigTh{n}$
expected time on every input, and $\bigTh{n^{2}}$ in the worst case.
\end{theorem}

\begin{proof}
\textbf{Worst case.} If every partition is maximally unbalanced, the range
shrinks by one each iteration and the partitions cost
$n, n-1, n-2, \ldots$, summing to $\bigTh{n^{2}}$.

\smallskip
\textbf{Expected case.} Call a pivot \emph{good} if it falls in the middle
half of the range, that is between the 25th and 75th percentiles. A uniformly
random pivot is good with probability $1/2$, and a good pivot leaves a
subproblem of size at most $\tfrac34$ of the range. Hence the expected number
of partitions needed to reduce the range by a factor of $\tfrac34$ is at most
2, and if $T(n)$ denotes the expected cost,
\[ T(n) \;\le\; T\bigl(\tfrac34 n\bigr) + \bigTh{n}, \]
where the $\bigTh{n}$ absorbs the expected two partitions. By
\chref{master} with $a = 1$, $b = 4/3$, $f(n) = n$: the watershed is
$n^{\log_{4/3} 1} = n^{0} = 1$, and $f(n) = n = \bigOm{n^{0+\varepsilon}}$ with
$\varepsilon = 1$, so case 3 applies provided regularity holds ---
$1 \cdot f(3n/4) = \tfrac34 n \le c\,n$ with $c = \tfrac34 < 1$, which it does.
Therefore $T(n) = \bigTh{f(n)} = \bigTh{n}$.
\end{proof}

\begin{conceptbox}[One deleted line, one changed case]
This is worth stating plainly, because it is the whole idea.

\smallskip
\textbf{Quicksort} solves both subproblems: $T(n) = 2T(n/2) + n$. Watershed
$n^{\log_{2}2} = n$, equal to $f(n)$, so \textbf{case 2} --- and case 2 is the
one that produces a logarithm. Result $\bigTh{n\lg n}$.

\smallskip
\textbf{Quickselect} solves one: $T(n) = T(n/2) + n$. Watershed
$n^{\log_{2}1} = 1$, far below $f(n) = n$, so \textbf{case 3} --- the root
dominates and the sum is geometric. Result $\bigTh{n}$.

\smallskip
The $\lg n$ factor is not saved by cleverness. It is the number of levels in
the tree, and discarding one branch collapses the tree to a path whose work
halves each step: $n + n/2 + n/4 + \cdots < 2n$.
\end{conceptbox}

\section{Measured}

\begin{worked}[Comparisons to Find the Median]
\texttt{code/ch11\_select.cpp}, mean of 25 independent runs --- an
\emph{expected} value is being measured, and a single run of a randomised
algorithm is not an answer.

\rowcolors{2}{white}{lightbg}
\begin{center}
\begin{tabular}{@{}rrrrr@{}}
\toprule
\rowcolor{primary!15}
$n$ & comparisons & $/n$ & spread over runs & $/(n\lg n)$ \\
\midrule
   10{,}000 &      33{,}640 & 3.36 & 1.80--5.84 & 0.2532 \\
  100{,}000 &     318{,}800 & 3.19 & 1.78--5.09 & 0.1919 \\
1{,}000{,}000 &   3{,}344{,}844 & 3.34 & 2.01--5.67 & 0.1678 \\
10{,}000{,}000 &  32{,}623{,}584 & 3.26 & 1.91--5.60 & 0.1403 \\
\bottomrule
\end{tabular}
\end{center}

\smallskip
\textbf{The cost per element is flat at about 3.3} across three orders of
magnitude of $n$. That flatness \emph{is} linearity: $\bigTh{n}$ means
comparisons$/n$ is bounded, and a column that neither rises nor falls over a
thousandfold range of $n$ is the evidence. The classical expected value for
median selection is $3.39n$.

\smallskip
\textbf{Meanwhile $/(n\lg n)$ falls steadily} --- 0.2532 to 0.1403 --- which is
what a linear quantity does when divided by an $n\lg n$ one.

\smallskip
\textbf{Read the spread column too.} Individual runs ranged from $1.8n$ to
$5.8n$ comparisons. The expectation is 3.3; the variance is real, and a single
unlucky run costs three times a lucky one. Expected-time guarantees are about
the average of many runs, not about the next one.
\end{worked}

\begin{worked}[Selecting Against Sorting, in Seconds]
\rowcolors{2}{white}{lightbg}
\begin{center}
\begin{tabular}{@{}rrrrr@{}}
\toprule
\rowcolor{primary!15}
$n$ & quickselect & \texttt{nth\_element} & full sort & speed-up \\
\midrule
  100{,}000 &  0.99 &  0.95 &   6.79 & 6.84 \\
  800{,}000 &  7.76 &  6.12 &  61.26 & 7.89 \\
3{,}200{,}000 & 29.94 & 29.76 & 284.57 & 9.50 \\
6{,}400{,}000 & 62.83 & 64.42 & 584.80 & 9.31 \\
\bottomrule
\end{tabular}
\end{center}

\smallskip
\textbf{Selection beats sorting by 7--9.5$\times$, and the factor grows} with
$n$, because it is essentially $\lg n$ divided by a constant.

\smallskip
\textbf{And the hand-written quickselect matches \texttt{std::nth\_element}}
--- within a few per cent either way. The library version is introselect, which
adds a median-of-medians fallback for the worst case; on random data that
fallback never fires, so the two are doing the same work.
\end{worked}

\begin{worked}[The Worst Case, and What It Costs]
A last-element pivot on already-sorted input, selecting the median:

\rowcolors{2}{white}{lightbg}
\begin{center}
\begin{tabular}{@{}rrrrr@{}}
\toprule
\rowcolor{primary!15}
$n$ & naive comparisons & $/(n^{2}/2)$ & random pivot & $/n$ \\
\midrule
 2{,}000 &      1{,}499{,}500 & 0.7498 &   9{,}402 & 4.70 \\
 8{,}000 &     23{,}998{,}000 & 0.7499 &  18{,}914 & 2.36 \\
16{,}000 &     95{,}996{,}000 & 0.7500 &  67{,}744 & 4.23 \\
32{,}000 &    383{,}992{,}000 & 0.7500 & 137{,}461 & 4.30 \\
\bottomrule
\end{tabular}
\end{center}

\smallskip
\textbf{The naive count is $0.7500 \times n^{2}/2 = 3n^{2}/8$ exactly}, and
that number can be derived. On sorted input the last-element pivot is always
the largest, so the range shrinks from the top by one each time. To reach the
median the range must fall from $n$ to $n/2$, costing
\[ n + (n-1) + \cdots + \frac{n}{2} \;=\; \frac{(n + n/2)}{2}\cdot\frac{n}{2}
   \;=\; \frac{3n^{2}}{8}. \]
Measured $0.7500$, predicted $0.75$. The analysis is not approximately right;
it is right.

\smallskip
\textbf{One difference from \chref{quicksort} worth noting.} There the same
input caused a \emph{crash}, because the recursion was $n$ deep. Here
\algname{Quickselect} is written as a loop, so the worst case is merely slow.
Writing a tail call as a loop did not fix the algorithm --- it is still
quadratic --- but it converted a crash into a performance problem, which is a
real improvement in its own right.
\end{worked}

\begin{worked}[Does $k$ Matter?]
Mean of 25 runs, $n = 1{,}000{,}000$:

\rowcolors{2}{white}{lightbg}
\begin{center}
\begin{tabular}{@{}rrr@{}}
\toprule
\rowcolor{primary!15}
$k/n$ & comparisons & $/n$ \\
\midrule
0.00 & 2{,}124{,}066 & 2.12 \\
0.10 & 2{,}615{,}691 & 2.62 \\
0.25 & 3{,}622{,}715 & 3.62 \\
0.50 & 2{,}992{,}344 & 2.99 \\
0.99 & 1{,}911{,}099 & 1.91 \\
1.00 & 1{,}945{,}088 & 1.95 \\
\bottomrule
\end{tabular}
\end{center}

\smallskip
\textbf{The ends are cheaper than the middle} --- about $2n$ for the minimum or
maximum against about $3n$ for the median --- but not dramatically so, and the
variation is less than a factor of two.

\smallskip
\textbf{The reason is a floor nothing can remove:} the \emph{first} partition
examines all $n$ elements whatever $k$ is. Even finding the minimum by
quickselect costs at least $n$ comparisons in that first pass, and the
remaining passes add about another $n$. So the answer to ``is it cheaper to
find the 10th smallest than the median?'' is yes, by about a third --- not by
the factor people expect.

\smallskip
(For the minimum specifically, a single sweep costs exactly $n-1$ comparisons
and is the obvious choice. Quickselect is for the cases a sweep cannot do.)
\end{worked}

\section{Linear in the Worst Case}

Quickselect's expected bound is excellent and its worst case is quadratic. A
deterministic algorithm achieves $\bigTh{n}$ in the \emph{worst} case.

\begin{conceptbox}[Median of medians]
To partition well you need a good pivot. To find a good pivot you could find
the median --- but that is the problem you are trying to solve. The way out is
to find a pivot that is merely \emph{guaranteed decent}, recursively.

\begin{tightnum}
  \item Divide the $n$ elements into $\lceil n/5 \rceil$ groups of 5.
  \item Find each group's median by sorting 5 elements --- $\bigO{1}$ each,
        $\bigO{n}$ in total.
  \item \textbf{Recursively select} the median of those $\lceil n/5 \rceil$
        medians. Call it $x$.
  \item Partition around $x$ and recurse into the side containing the answer.
\end{tightnum}

\smallskip
\textbf{Why $x$ is a decent pivot.} At least half the groups have their median
$\le x$, and in each such group at least 3 of the 5 elements are $\le$ that
median. So at least $3 \cdot \tfrac12 \cdot \tfrac{n}{5} = \tfrac{3n}{10}$
elements are $\le x$, and symmetrically $\tfrac{3n}{10}$ are $\ge x$.
Therefore the recursive call in step 4 is on at most $\tfrac{7n}{10}$
elements --- always, with no probability involved.

\smallskip
\textbf{The recurrence} is then
\[ T(n) \;\le\; \underbrace{T(n/5)}_{\text{step 3}}
             + \underbrace{T(7n/10)}_{\text{step 4}}
             + \bigTh{n}, \]
which the master theorem cannot touch (two different subproblem sizes). But
$\tfrac15 + \tfrac{7}{10} = \tfrac{9}{10} < 1$, and substitution with the guess
$T(n) \le cn$ gives $T(n) \le \tfrac{c n}{5} + \tfrac{7cn}{10} + an =
\tfrac{9cn}{10} + an \le cn$ whenever $c \ge 10a$. So $T(n) = \bigO{n}$.

\smallskip
\textbf{The groups of 5 are not arbitrary.} With groups of 3 the sum becomes
$\tfrac13 + \tfrac23 = 1$, the substitution fails, and the bound is
$\bigTh{n\lg n}$. The constant 5 is the smallest odd group size that makes the
fractions sum to less than 1.
\end{conceptbox}

\begin{alertbox}[And yet almost nobody uses it]
Median-of-medians is $\bigTh{n}$ worst case, which is strictly better than
quickselect's guarantee. Its constant is roughly 5--10$\times$ larger: it sorts
every group of five, makes a second recursive call, and touches memory in a
scattered pattern.

\smallskip
So in practice:
\begin{tight}
  \item \textbf{randomised quickselect} for ordinary work, at about $3.3n$
        comparisons measured;
  \item \textbf{introselect} --- quickselect that counts its iterations and
        switches to median-of-medians if progress is too slow --- for library
        code. This is what \texttt{std::nth\_element} does, and it is the same
        design as \texttt{std::sort}'s heapsort fallback in \chref{quicksort}.
\end{tight}

\smallskip
The pattern is worth naming, because it recurs: \textbf{use the algorithm with
the better constant, and keep the one with the better guarantee as a
fallback.} You get the common case fast and the worst case bounded.
\end{alertbox}

\begin{pitfall}
\begin{tight}
  \item \textbf{Sorting to find one element.} Measured, 7--9.5$\times$ more
        expensive, growing with $n$.
  \item \textbf{Recursing into both sides.} That is quicksort. The deleted call
        is the entire algorithm.
  \item \textbf{Judging an expected bound from one run.} Measured runs ranged
        from $1.8n$ to $5.8n$ for a mean of $3.3n$.
  \item \textbf{Using a fixed pivot.} Sorted input then gives $3n^{2}/8$
        comparisons, measured exactly.
  \item \textbf{Reaching for median-of-medians by default.} Its constant is
        5--10$\times$ worse; it is a fallback, not a first choice.
  \item \textbf{Using groups of 3.} The fractions sum to 1 and the bound
        degrades to $\bigTh{n\lg n}$.
  \item \textbf{Expecting selection to be much cheaper near the ends.}
        Measured $2.1n$ at $k=0$ against $3.0n$ at the median --- the first
        partition costs $n$ regardless.
\end{tight}
\end{pitfall}

\begin{lab}[Selection Against Sorting]
Reproduces all four tables. Expect twenty-five minutes.

\begin{lstlisting}[language=C++]
// ch11_select.cpp -- finding the k-th smallest without sorting.
// Build: cl /O2 /EHsc /std:c++17 ch11_select.cpp     (see Appendix A)

// --- 1. Partition, exactly as in Chapter 9 -------------------------

// --- 2. Quickselect: recurse into ONE side, not both ---------------
// This single change is the whole algorithm. Quicksort solves both
// subproblems and gets T(n) = 2T(n/2) + n = Theta(n lg n).
// Quickselect discards one and gets T(n) = T(n/2) + n, which the
// master theorem settles as Theta(n) -- case 3, the root dominating.
//
// Written as a loop rather than a recursion, because the recursive
// call is in tail position and there is no reason to spend a stack
// frame on it. That also removes Chapter 9's stack-overflow failure.
static int quickselect(std::vector<int>& a, size_t k, bool randomised) {
    size_t lo = 0, hi = a.size();
    for (;;) {
        if (hi - lo == 1) return a[lo];
        size_t p = randomised ? partition_random(a, lo, hi)
                              : partition_last(a, lo, hi);
        if (p == k) return a[p];
        if (k < p) hi = p;          // the k-th is on the left
        else       lo = p + 1;      // ... or on the right
    }
}

// --- 3. Comparisons: selection against sorting ---------------------
// An EXPECTED value is being measured, so one run of a randomised
// algorithm is not an answer: single runs of this gave 2.52, 3.92,
// 3.73 and 2.73 comparisons per element, which says only that the
// answer is somewhere near 3. Average over 25 independent runs.

// --- 4. The same thing in seconds ----------------------------------
// std::nth_element is the library's version: introselect, which is
// quickselect with a median-of-medians fallback, exactly as
// std::sort is quicksort with a heapsort fallback.

// --- 5. The worst case, for Chapter 9's reason ---------------------
// A last-element pivot on sorted input shrinks the range by one each
// time: T(n) = T(n-1) + n = Theta(n^2). The loop form means it is
// slow rather than a crash, which is an improvement but not a fix.

// --- 6. Does k matter? ---------------------------------------------
// The expected cost is smaller towards the ends -- but not
// dramatically, because the first partition costs n whatever k is.
// That first n is a floor nothing removes.
\end{lstlisting}

\textbf{What to notice.} The spread column in part 3. The mean is a stable 3.3
across four orders of magnitude, and individual runs vary by a factor of three.
Both facts are properties of the algorithm, and quoting only the first would be
misleading.

\smallskip
\textbf{Then try to break it.} Implement median-of-medians and race it against
randomised quickselect on random input. Predict first which wins and by how
much. Then build the input that makes quickselect quadratic and race them
again. The crossover between ``better constant'' and ``better guarantee'' is the
entire argument for introselect, and you will have measured both ends of it.
\end{lab}

\begin{checkpoint}
\begin{tightnum}
  \item Quicksort's recurrence is $T(n)=2T(n/2)+n$ and quickselect's is
        $T(n)=T(n/2)+n$. Identify the master-theorem case for each and explain
        where the $\lg n$ went.
  \item Derive the $3n^{2}/8$ comparison count for a last-element pivot
        selecting the median of sorted input.
  \item Why does median-of-medians use groups of 5 rather than 3? Show the
        arithmetic that fails.
\end{tightnum}
\end{checkpoint}

\begin{chaptersummary}
\begin{tight}
  \item The $i$-th order statistic is the $i$-th smallest. Selection asks for
        one of them; sorting computes all of them and is asymptotically more
        expensive.
  \item Quickselect is quicksort with one recursive call deleted. That moves
        the recurrence from master-theorem case 2 ($\bigTh{n\lg n}$) to case 3
        ($\bigTh{n}$): the tree collapses to a path and the work is geometric.
  \item Expected $\bigTh{n}$ on every input with a random pivot;
        $\bigTh{n^{2}}$ worst case.
  \item Measured: $3.3$ comparisons per element for the median, flat across
        $n = 10^{4}$ to $10^{7}$ --- and individual runs ranged from $1.8n$ to
        $5.8n$.
  \item Measured: 7--9.5$\times$ faster than a full sort, a factor that grows
        with $n$; and indistinguishable from \texttt{std::nth\_element}.
  \item Measured: the naive pivot on sorted input used exactly $3n^{2}/8$
        comparisons, matching the derivation to four decimal places.
  \item Median-of-medians is worst-case $\bigTh{n}$ via $T(n) \le T(n/5) +
        T(7n/10) + \bigTh{n}$, where $\tfrac15+\tfrac{7}{10}<1$ is what makes
        the substitution close. Its constant is 5--10$\times$ worse, so it is
        used as a fallback, not a default.
  \item Cost varies with $k$ by less than a factor of two, because the first
        partition costs $n$ whatever $k$ is.
\end{tight}
\end{chaptersummary}

\begin{reviewq}
  \item \textbf{(MCQ)} Randomised quickselect's expected running time is:
        (a)~$\bigTh{\lg n}$; (b)~$\bigTh{n}$; (c)~$\bigTh{n\lg n}$;
        (d)~$\bigTh{n^{2}}$.
  \item \textbf{(MCQ)} Median-of-medians partitions around an element
        guaranteed to exceed at least: (a)~$n/5$; (b)~$3n/10$; (c)~$n/2$;
        (d)~$7n/10$ of the elements.
  \item \textbf{(Short)} Explain why deleting one recursive call changes the
        master-theorem case, and hence removes the logarithm.
  \item \textbf{(Short)} State the trade-off between randomised quickselect and
        median-of-medians, and say what introselect does about it.
  \item \textbf{(Short)} The measured spread was $1.8n$ to $5.8n$ for a mean of
        $3.3n$. What does that tell you about relying on an expected-time
        bound for a single latency-critical request?
  \item \textbf{(Applied)} Prove that at least $n-1$ comparisons are needed to
        find the minimum of $n$ elements, and explain why you would use a
        single sweep rather than quickselect for $k = 1$.
  \item \textbf{(Applied)} You need the 99th percentile latency from 50 million
        recorded request times, once per minute. Choose an algorithm, justify
        it with the measurements in this chapter, and say what would change
        your choice.
\end{reviewq}
